\IfFileExists{stacks-project.cls}{%
\documentclass{stacks-project}
}{%
\documentclass{amsart}
}

% For dealing with references we use the comment environment
\usepackage{verbatim}
\newenvironment{reference}{\comment}{\endcomment}
%\newenvironment{reference}{}{}
\newenvironment{slogan}{\comment}{\endcomment}
\newenvironment{history}{\comment}{\endcomment}

% For commutative diagrams we use Xy-pic
\usepackage[all]{xy}

% We use 2cell for 2-commutative diagrams.
\xyoption{2cell}
\UseAllTwocells

% We use multicol for the list of chapters between chapters
\usepackage{multicol}

% This is generally recommended for better output
\usepackage{lmodern}
\usepackage[T1]{fontenc}

% For cross-file-references
\usepackage{xr-hyper}

% Package for hypertext links:
\usepackage{hyperref}

% For any local file, say "hello.tex" you want to link to please
% use \externaldocument[hello-]{hello}
\externaldocument[introduction-]{introduction}
\externaldocument[conventions-]{conventions}
\externaldocument[sets-]{sets}
\externaldocument[categories-]{categories}
\externaldocument[topology-]{topology}
\externaldocument[sheaves-]{sheaves}
\externaldocument[sites-]{sites}
\externaldocument[stacks-]{stacks}
\externaldocument[fields-]{fields}
\externaldocument[algebra-]{algebra}
\externaldocument[brauer-]{brauer}
\externaldocument[homology-]{homology}
\externaldocument[derived-]{derived}
\externaldocument[simplicial-]{simplicial}
\externaldocument[more-algebra-]{more-algebra}
\externaldocument[smoothing-]{smoothing}
\externaldocument[modules-]{modules}
\externaldocument[sites-modules-]{sites-modules}
\externaldocument[injectives-]{injectives}
\externaldocument[cohomology-]{cohomology}
\externaldocument[sites-cohomology-]{sites-cohomology}
\externaldocument[dga-]{dga}
\externaldocument[dpa-]{dpa}
\externaldocument[sdga-]{sdga}
\externaldocument[hypercovering-]{hypercovering}
\externaldocument[schemes-]{schemes}
\externaldocument[constructions-]{constructions}
\externaldocument[properties-]{properties}
\externaldocument[morphisms-]{morphisms}
\externaldocument[coherent-]{coherent}
\externaldocument[divisors-]{divisors}
\externaldocument[limits-]{limits}
\externaldocument[varieties-]{varieties}
\externaldocument[topologies-]{topologies}
\externaldocument[descent-]{descent}
\externaldocument[perfect-]{perfect}
\externaldocument[more-morphisms-]{more-morphisms}
\externaldocument[flat-]{flat}
\externaldocument[groupoids-]{groupoids}
\externaldocument[more-groupoids-]{more-groupoids}
\externaldocument[etale-]{etale}
\externaldocument[chow-]{chow}
\externaldocument[intersection-]{intersection}
\externaldocument[pic-]{pic}
\externaldocument[weil-]{weil}
\externaldocument[adequate-]{adequate}
\externaldocument[dualizing-]{dualizing}
\externaldocument[duality-]{duality}
\externaldocument[discriminant-]{discriminant}
\externaldocument[derham-]{derham}
\externaldocument[local-cohomology-]{local-cohomology}
\externaldocument[algebraization-]{algebraization}
\externaldocument[curves-]{curves}
\externaldocument[resolve-]{resolve}
\externaldocument[models-]{models}
\externaldocument[functors-]{functors}
\externaldocument[equiv-]{equiv}
\externaldocument[pione-]{pione}
\externaldocument[etale-cohomology-]{etale-cohomology}
\externaldocument[proetale-]{proetale}
\externaldocument[relative-cycles-]{relative-cycles}
\externaldocument[more-etale-]{more-etale}
\externaldocument[trace-]{trace}
\externaldocument[crystalline-]{crystalline}
\externaldocument[spaces-]{spaces}
\externaldocument[spaces-properties-]{spaces-properties}
\externaldocument[spaces-morphisms-]{spaces-morphisms}
\externaldocument[decent-spaces-]{decent-spaces}
\externaldocument[spaces-cohomology-]{spaces-cohomology}
\externaldocument[spaces-limits-]{spaces-limits}
\externaldocument[spaces-divisors-]{spaces-divisors}
\externaldocument[spaces-over-fields-]{spaces-over-fields}
\externaldocument[spaces-topologies-]{spaces-topologies}
\externaldocument[spaces-descent-]{spaces-descent}
\externaldocument[spaces-perfect-]{spaces-perfect}
\externaldocument[spaces-more-morphisms-]{spaces-more-morphisms}
\externaldocument[spaces-flat-]{spaces-flat}
\externaldocument[spaces-groupoids-]{spaces-groupoids}
\externaldocument[spaces-more-groupoids-]{spaces-more-groupoids}
\externaldocument[bootstrap-]{bootstrap}
\externaldocument[spaces-pushouts-]{spaces-pushouts}
\externaldocument[spaces-chow-]{spaces-chow}
\externaldocument[groupoids-quotients-]{groupoids-quotients}
\externaldocument[spaces-more-cohomology-]{spaces-more-cohomology}
\externaldocument[spaces-simplicial-]{spaces-simplicial}
\externaldocument[spaces-duality-]{spaces-duality}
\externaldocument[formal-spaces-]{formal-spaces}
\externaldocument[restricted-]{restricted}
\externaldocument[spaces-resolve-]{spaces-resolve}
\externaldocument[formal-defos-]{formal-defos}
\externaldocument[defos-]{defos}
\externaldocument[cotangent-]{cotangent}
\externaldocument[examples-defos-]{examples-defos}
\externaldocument[algebraic-]{algebraic}
\externaldocument[examples-stacks-]{examples-stacks}
\externaldocument[stacks-sheaves-]{stacks-sheaves}
\externaldocument[criteria-]{criteria}
\externaldocument[artin-]{artin}
\externaldocument[quot-]{quot}
\externaldocument[stacks-properties-]{stacks-properties}
\externaldocument[stacks-morphisms-]{stacks-morphisms}
\externaldocument[stacks-limits-]{stacks-limits}
\externaldocument[stacks-cohomology-]{stacks-cohomology}
\externaldocument[stacks-perfect-]{stacks-perfect}
\externaldocument[stacks-introduction-]{stacks-introduction}
\externaldocument[stacks-more-morphisms-]{stacks-more-morphisms}
\externaldocument[stacks-geometry-]{stacks-geometry}
\externaldocument[moduli-]{moduli}
\externaldocument[moduli-curves-]{moduli-curves}
\externaldocument[examples-]{examples}
\externaldocument[exercises-]{exercises}
\externaldocument[guide-]{guide}
\externaldocument[desirables-]{desirables}
\externaldocument[coding-]{coding}
\externaldocument[obsolete-]{obsolete}
\externaldocument[fdl-]{fdl}
\externaldocument[index-]{index}

% Theorem environments.
%
\theoremstyle{plain}
\newtheorem{theorem}[subsection]{Theorem}
\newtheorem{proposition}[subsection]{Proposition}
\newtheorem{lemma}[subsection]{Lemma}

\theoremstyle{definition}
\newtheorem{definition}[subsection]{Definition}
\newtheorem{example}[subsection]{Example}
\newtheorem{exercise}[subsection]{Exercise}
\newtheorem{situation}[subsection]{Situation}

\theoremstyle{remark}
\newtheorem{remark}[subsection]{Remark}
\newtheorem{remarks}[subsection]{Remarks}

\numberwithin{equation}{subsection}

% Macros
%
\def\lim{\mathop{\mathrm{lim}}\nolimits}
\def\colim{\mathop{\mathrm{colim}}\nolimits}
\def\Spec{\mathop{\mathrm{Spec}}}
\def\Hom{\mathop{\mathrm{Hom}}\nolimits}
\def\Ext{\mathop{\mathrm{Ext}}\nolimits}
\def\SheafHom{\mathop{\mathcal{H}\!\mathit{om}}\nolimits}
\def\SheafExt{\mathop{\mathcal{E}\!\mathit{xt}}\nolimits}
\def\Sch{\mathit{Sch}}
\def\Mor{\mathop{\mathrm{Mor}}\nolimits}
\def\Ob{\mathop{\mathrm{Ob}}\nolimits}
\def\Sh{\mathop{\mathit{Sh}}\nolimits}
\def\NL{\mathop{N\!L}\nolimits}
\def\CH{\mathop{\mathrm{CH}}\nolimits}
\def\proetale{{pro\text{-}\acute{e}tale}}
\def\etale{{\acute{e}tale}}
\def\QCoh{\mathit{QCoh}}
\def\Ker{\mathop{\mathrm{Ker}}}
\def\Im{\mathop{\mathrm{Im}}}
\def\Coker{\mathop{\mathrm{Coker}}}
\def\Coim{\mathop{\mathrm{Coim}}}

% Boxtimes
%
\DeclareMathSymbol{\boxtimes}{\mathbin}{AMSa}{"02}

%
% Macros for moduli stacks/spaces
%
\def\QCohstack{\mathcal{QC}\!\mathit{oh}}
\def\Cohstack{\mathcal{C}\!\mathit{oh}}
\def\Spacesstack{\mathcal{S}\!\mathit{paces}}
\def\Quotfunctor{\mathrm{Quot}}
\def\Hilbfunctor{\mathrm{Hilb}}
\def\Curvesstack{\mathcal{C}\!\mathit{urves}}
\def\Polarizedstack{\mathcal{P}\!\mathit{olarized}}
\def\Complexesstack{\mathcal{C}\!\mathit{omplexes}}
% \Pic is the operator that assigns to X its picard group, usage \Pic(X)
% \Picardstack_{X/B} denotes the Picard stack of X over B
% \Picardfunctor_{X/B} denotes the Picard functor of X over B
\def\Pic{\mathop{\mathrm{Pic}}\nolimits}
\def\Picardstack{\mathcal{P}\!\mathit{ic}}
\def\Picardfunctor{\mathrm{Pic}}
\def\Deformationcategory{\mathcal{D}\!\mathit{ef}}
\usepackage{cleveref}
\crefname{originalreference}{original reference}{original references}
\Crefname{originalreference}{Original reference}{Original references}
\makeatletter
\newlabel{lemma-exactness-g-shriek-modules}{{0D74}{source}{}{sourcepointer-1}{}}
\newlabel{lemma-exactness-g-shriek-modules@cref}{{[originalreference][0][]0D74}{[1][0][]source}{}{}{}}
\makeatother
\crefname{originalreference}{original reference}{original references}
\Crefname{originalreference}{Original reference}{Original references}
\makeatletter
\newlabel{lemma-restriction-module-to-components-site}{{0D72}{source}{}{sourcepointer-2}{}}
\newlabel{lemma-restriction-module-to-components-site@cref}{{[originalreference][0][]0D72}{[1][0][]source}{}{}{}}
\makeatother
\crefname{originalreference}{original reference}{original references}
\Crefname{originalreference}{Original reference}{Original references}
\makeatletter
\newlabel{lemma-restriction-to-components-site}{{09WG}{source}{}{sourcepointer-3}{}}
\newlabel{lemma-restriction-to-components-site@cref}{{[originalreference][0][]09WG}{[1][0][]source}{}{}{}}
\makeatother
\crefname{originalreference}{original reference}{original references}
\Crefname{originalreference}{Original reference}{Original references}
\makeatletter
\newlabel{situation-simplicial-site}{{09WE}{source}{}{sourcepointer-4}{}}
\newlabel{situation-simplicial-site@cref}{{[originalreference][0][]09WE}{[1][0][]source}{}{}{}}
\makeatother
\title{Cartesian objects of the derived category characterized by the shriek-tensor maps g\_\{n!\}K\_n -> g\_\{n!\}Z tensor K}
\author{The Stacks Project Authors}
\date{stacks-project:spaces-simplicial.tex:local-snapshot-20261012; GNU FDL 1.2 or later}
\begin{document}
\maketitle
\noindent\textit{Editorial note.} All delivered mathematical text is the author's own, copied from the pinned Stacks Project chapter file sources/stacks-spaces-simplicial/pinned-original-source.tex; no statement or proof step was rewritten, reordered or removed.\par\medskip
\noindent\textit{Editorial note.} The counted proof is the author's complete argument: the final equalities in the displayed restrictions (f\_varphi\textasciicircum{}*O\_n tensor\_\{O\_m\} K\_m = K\_m, source lines 2806-2809) use that for a morphism of ringed topoi the pullback of the structure sheaf is the target structure sheaf (f\_varphi\textasciicircum{}*O\_n = O\_m), as set up in the lemma statement; the 'Thus the statement is clear' lines close arguments whose entire componentwise content is displayed immediately above; part (2) additionally needs the flatness assumption to keep the derived tensor termwise on components, as stated.\par\medskip
\noindent\textit{Editorial note.} The label lemma-derived-cartesian-shriek names the counted statement, and the proof environment immediately following it in the source is delivered verbatim as its proof.\par\medskip
\noindent\textit{Editorial note.} The counted statement is carried by the author's own lemma environment in the chapter source.\par\medskip
\noindent\textit{Allowed standard prerequisites (editorial).}\begin{itemize}
\item Exactness of the component shrieks. Same chapter, \textbackslash{}label\{lemma-exactness-g-shriek-modules\} and the exactness part of \textbackslash{}label\{lemma-restriction-to-components-site\} (card L46149). Lines 2776-2778 and 2791-2794: g\_\{n!\} exact induces a derived functor adjoint to g\_n\textasciicircum{}\{-1\} resp. g\_n\textasciicircum{}*.
\item The component formula for g\_\{n!\}. Same chapter, \textbackslash{}label\{lemma-restriction-to-components-site\} (card L46149) and \textbackslash{}label\{lemma-restriction-module-to-components-site\} (card L46168 of this batch). Lines 2782-2787 and 2799-2809: the restriction to C\_m of g\_\{n!\}K\_n and of the tensor product is computed by the direct sums over varphi : [n] -> [m].
\item Derived tensor product and its restriction to components. Standard interaction of derived tensor with exact functors and restrictions. Lines 2799-2809: the restriction of g\_\{n!\}O\_n tensor\textasciicircum{}L\_O K to C\_m is the termwise tensor over O\_m.
\item The unit map of the (g\_\{n!\}, g\_n\textasciicircum{}\{-1\}) adjunction. Standard adjunction unit for the shriek-pushforward pair. Lines 2778-2782: the canonical map is the adjoint of Z -> g\_n\textasciicircum{}\{-1\}g\_\{n!\}Z adjoint to id\_\{g\_\{n!\}Z\}.
\end{itemize}\medskip
\ifdefined\hypertarget\hypertarget{sourcepointer-1}{}\fi
\noindent\textit{Source reference.} Original reference 0D74: sources/stacks-spaces-simplicial/pinned-original-source.tex, source lines 1533--1533. Prerequisite pointer invoked at proof line(s) 2795; it names a standard fact of the same chapter that the counted argument uses, and every step of the counted proof built on it is retained verbatim in the delivered body.\par\medskip
\ifdefined\hypertarget\hypertarget{sourcepointer-2}{}\fi
\noindent\textit{Source reference.} Original reference 0D72: sources/stacks-spaces-simplicial/pinned-original-source.tex, source lines 1459--1459. Prerequisite pointer invoked at proof line(s) 2803; it names a standard fact of the same chapter that the counted argument uses, and every step of the counted proof built on it is retained verbatim in the delivered body.\par\medskip
\ifdefined\hypertarget\hypertarget{sourcepointer-3}{}\fi
\noindent\textit{Source reference.} Original reference 09WG: sources/stacks-spaces-simplicial/pinned-original-source.tex, source lines 970--970. Prerequisite pointer invoked at proof line(s) 2783; it names a standard fact of the same chapter that the counted argument uses, and every step of the counted proof built on it is retained verbatim in the delivered body.\par\medskip
\ifdefined\hypertarget\hypertarget{sourcepointer-4}{}\fi
\noindent\textit{Source reference.} Original reference 09WE: sources/stacks-spaces-simplicial/pinned-original-source.tex, source lines 888--888. Prerequisite pointer invoked at proof line(s) 2753; it names a standard fact of the same chapter that the counted argument uses, and every step of the counted proof built on it is retained verbatim in the delivered body.\par\medskip
\begin{lemma}
\label{lemma-derived-cartesian-shriek}
In Situation \ref{situation-simplicial-site}.
\begin{enumerate}
\item An object $K$ of $D(\mathcal{C}_{total})$ is cartesian if and only
the canonical map
$$
g_{n!}K_n \longrightarrow
g_{n!}\mathbf{Z} \otimes^\mathbf{L}_\mathbf{Z} K
$$
is an isomorphism for all $n$.
\item Let $\mathcal{O}$ be a sheaf of rings on $\mathcal{C}_{total}$
such that the morphisms $f_\varphi^{-1}\mathcal{O}_n \to \mathcal{O}_m$
are flat for all $\varphi : [n] \to [m]$. Then an object $K$ of
$D(\mathcal{O})$ is cartesian if and only if the canonical map
$$
g_{n!}K_n \longrightarrow
g_{n!}\mathcal{O}_n \otimes^\mathbf{L}_\mathcal{O} K
$$
is an isomorphism for all $n$.
\end{enumerate}
\end{lemma}

\begin{proof}
Proof of (1). Since $g_{n!}$ is exact, it induces a functor
on derived categories adjoint to $g_n^{-1}$.
The map is the adjoint of the map
$K_n \to (g_n^{-1}g_{n!}\mathbf{Z}) \otimes^\mathbf{L}_\mathbf{Z} K_n$
corresponding to $\mathbf{Z} \to g_n^{-1}g_{n!}\mathbf{Z}$
which in turn is adjoint to
$\text{id} : g_{n!}\mathbf{Z} \to g_{n!}\mathbf{Z}$.
Using the description of $g_{n!}$
given in Lemma \ref{lemma-restriction-to-components-site}
we see that the restriction to $\mathcal{C}_m$ of this map
is
$$
\bigoplus\nolimits_{\varphi : [n] \to [m]} f_\varphi^{-1}K_n
\longrightarrow
\bigoplus\nolimits_{\varphi : [n] \to [m]} K_m
$$
Thus the statement is clear.

\medskip\noindent
Proof of (2). Since $g_{n!}$ is exact
(Lemma \ref{lemma-exactness-g-shriek-modules}), it induces a functor
on derived categories adjoint to $g_n^*$ (also exact).
The map is the adjoint of the map
$K_n \to (g_n^*g_{n!}\mathcal{O}_n) \otimes^\mathbf{L}_{\mathcal{O}_n} K_n$
corresponding to $\mathcal{O}_n \to g_n^*g_{n!}\mathcal{O}_n$
which in turn is adjoint to
$\text{id} : g_{n!}\mathcal{O}_n \to g_{n!}\mathcal{O}_n$.
Using the description of $g_{n!}$
given in Lemma \ref{lemma-restriction-module-to-components-site}
we see that the restriction to $\mathcal{C}_m$ of this map
is
$$
\bigoplus\nolimits_{\varphi : [n] \to [m]} f_\varphi^*K_n
\longrightarrow
\bigoplus\nolimits_{\varphi : [n] \to [m]}
f_\varphi^*\mathcal{O}_n \otimes_{\mathcal{O}_m} K_m =
\bigoplus\nolimits_{\varphi : [n] \to [m]} K_m
$$
Thus the statement is clear.
\end{proof}

\end{document}
